\documentstyle[margin,line]{resume}

% Logos
% =====

\font\attand=cmr7
%\font\MetafontLogoFont=logo10

\def\PS{{\tt P\small OST\tt S\small CRIPT}}
\def\ATT{{AT{\attand \&}T}}
\def\MF{{METAFONT}}
\def\Cplusplus{{\rm C\raise.5ex\hbox{\small ++}}}
\def\AmSTeX{{$\cal A\kern-.1667em\lower.5ex\hbox{$\cal M$}\kern-.125em
S$-\TeX}}

\oddsidemargin -0.7in
\evensidemargin -0.7in
\textwidth=6.6in
\textheight = 10.2in

\begin{document}

\name{Tal Garfinkel}
\address{talg@cs.stanford.edu\\
(650) 575-6744\\
5681 North Verde Mountain Dr.\\
Tucson, AZ 85750}

\begin{resume}

\section{\sc Education}

{\bf Stanford University}
{\it Fall 2001 -- Spring 2010} \newline
{\it PhD Computer Science, 2010}

{\bf University of California, Berkeley }
{\it Fall 1998 -- Spring 2001 } \newline
{\it BA Computer Science with Honors, 2002}

%%%%%%%%%%%%% Work experience

\begin{format}
\title{l} \employer{r} \\ \location{l} \dates{r}\\
\body\\
\end{format}


\section{\sc Selected Employment History}

\title{\em Research Engineer}
\employer{\bf University of California, San Deigo}
\dates{Spring 2022--Spring 2025}
\location{San Diego, California}
\begin{position}
My research at UCSD has focuses on the development of new software and hardware
mechanisms for fine grain isolation, and other projects at the intersection of
systems and security.  \end{position}


\title{\em Principal Consultant}
\employer{\bf Corepoint Systems, LLC}
\dates{Winter 2016--2022}
\location{Palo Alto, California}
\begin{position}
Recent projects include assisting the VMware AppDefense team on the design of
their endpoint security product, work with the Mozilla
WebAssembly team on support for distributed programming in Wasm/WASI, and work with
colleagues at UCSD on library sandboxing.
\end{position}


\title{\em Research Associate}
\employer{\bf Stanford University}
\dates{Spring 2017--Fall 2019}
\location{Stanford, California}
\begin{position} I co-lead the development of the Castor, a low overhead
record/replay system for production multi-core applications, and supervised
graduate students and undergrads on the project. More recent projects
focused on improving memory safety in WebAssembly, and language and runtime
support for library sandboxing in Firefox.

\end{position}

\title{\em Expert Witness}
\employer{\bf Harbor Labs/Haynes and Boone, LLP}

\dates{Summer 2014 -- Spring 2015}
\location{Palo Alto, California}
\begin{position}
I worked as an expert witness on the patent owner side of an IPR proceeding
regarding network security, virtualization and replay.
\end{position}


\title{\em Software Engineer}
\employer{\bf Bebop}

\dates{Fall 2013 -- Spring 2014}
\location{Palo Alto, California}
\begin{position}
I worked as a full stack developer at a small startup doing end user focused
software for the Enterprise. I focused on authentication, session managment,
and other security and system level aspects of the product.
\end{position}


\title{\em Senior Software Engineer}
\employer{\bf VMWare Cloud Foundry Team}

\dates{Fall 2010 -- Fall 2012}
\location{Palo Alto, California}
\begin{position}
I worked as a developer on the Cloud Foundry ({\tt www.cloudfoundry.org}) core
team, developing an Open Source platform as a service. Our system was used
as the basis for our own PaaS offering, and is widely used in private cloud
infrastructure. My work spanned a range of areas including on app build caching
and deployment, data path offload, container managment, and security.

\end{position}


\title{\em Senior Scientist}
\employer{\bf VMWare Advanced Development}

\dates{Fall 06 -- Fall 08}
\location{Palo Alto, California}
\begin{position}
I developed, published, and patented new technologies, and consulted on
security architecture for the core platform, cloud computing and desktop
offerings.

\end{position}

\title{\em Graduate Research Assistant}
\employer{\bf Stanford University}

\dates{Summer 2001 -- Fall 2006, Fall 08 -- Spring 2010}
\location{Stanford, California}
\begin{position}
During my PhD, I was a research assistant working on topics in Computer Security
and Systems.  Advised by Prof. Mendel Rosenblum and Prof. Dan Boneh. I developed
new technologies for improving security in operating systems,
virtual machines, and networks.  \end{position}

\pagebreak
\title{\em Engineering Intern/Consultant}
\employer{\bf VMware Inc.}
\dates{Summer 2003--Fall 2006}
\location{Palo Alto, California}
\begin{position}
As an engineering intern and later consultant advised by Mendel Rosenblum I
worked on new product development and improving security in existing
VMware products.

\end{position}

%Brewer Research
\title{\em Undergraduate Research Assistant}
\employer{\bf University of California, Berkeley}
\dates{Spring 99 -- Spring 2001 }
\location{Berkeley, California}
\begin{position}
As an undergraduate research assistant advised by Prof. David Wagner, Prof.
Steven Gribble and Prof. Eric Brewer.  I lead the development of the current
version of Janus, a tool for confining network services in a restricted
execution environment.  I also developed a fast HTTP front end for a novel
asynchronous I/O system,  and researched mechanisms for fast user level
networking in clusters as well as other relevant design issues in building
scalable, high performance, Internet services.
\end{position}

% Debevec research

\title{\em Undergraduate Research Assistant}
\employer{\bf University of California, Berkeley}
\dates{Spring 98 -- Fall 99}
\location{Berkeley, California}
\begin{position}
As a research assistant in Graphics and Scalable systems, Advised by Dr.
Paul Debevec. I developed tools for distributed rendering, distributed
batch processing, remote execution and image processing on large
clusters. I also worked on the production of several animated short films
including significant contributions to Fiat Lux which has won numerous
awards and appeared in the SIGGRAPH 99 Electronic Theater.
\end{position}

%\title{\em Programmer/Consultant}
%\employer{\bf University of California, Berkeley }
%\dates{Spring 99}
%\location{Berkeley, California}
%\begin{position}
%I designed and developed a suite of small composable tools to analyze rainfall
%patterns for use in Hydrology research.
%\end{position}


%\title{\em NSF Research Internship }
%\employer{\bf Center for Quantized Electronic Structures}
%\dates{Summer 98}
%\location{University of California, Santa Barbara}
%\begin{position}
%As a research assistant in Computational Chemistry. Designed, implemented and
%analyzed a Monte Carlo model of iron dispersion
%in a ferromagnetic bromo-sodalite ${Na_8Br_2[AlSiO_4]_6}$.
%\end{position}

%\title{\em Software Engineer}
%\employer{\bf Utilicom Inc.}
%\dates{ Winter 96 -- Fall 97.}
%\location{Goleta, California}
%\begin{position}
%I developed embedded software for point to point wireless communication systems, PC based test software for wireless modems, and tools for finding optimal sets
%of codes for use in CDMA systems.
%\end{position}

%\title{\em Software Consultant}
%\employer{\bf Santa Barbara County Education Office}
%\dates{ Summer 95 -- Winter 96}
%\location{Santa Barbara, California}
%\begin{position}
%Developed relational databases for personnel data as well
%as performing various system administration tasks.
%\end{position}

\section{\sc Publications}

{\em Extended User Interrupts (xUI): Fast and Flexible Asynchronous
Notification without Polling.} Berk Aydogmus, Linsong Guo, Danial Zuberi, Tal
Garfinkel, Dean Tullsen, Amy Ousterhout, Kazem Taram.  In Architectural Support
for Programming Languages and Operating Systems (ASPLOS), 2025

{\em Segue \& ColorGuard: Optimizing SFI Performance and Scalability on Modern
Architectures.} Shravan Narayan, Tal Garfinkel, Evan Johnson, Zachary Yedidia,
Yingchen  Wang, Andrew Brown, Anjo Vahldiek-Oberwagner, Michael LeMay, Wenyong
Huang, Xin Wang, Mingqiu Sun,Dean Tullsen, Deian Stefan. In
Architectural Support for Programming Languages and Operating Systems (ASPLOS),
2025

{\em The Benefits and Limitations of User Interrupts for Preemptive Userspace
Scheduling.} Linsong Guo, Danial Zuberi, Tal Garfinkel, Amy Ousterhout.
In Network System Design and Implementation(NSDI), 2025

{\em Hardware-Assisted Fault Isolation: Going Beyond the Limits of
Software-Based Sandboxing.} Shravan Narayan, Tal Garfinkel, Mohammadkazem
Taram, Joey Rudek, Daniel Moghimi, Evan Johnson, Chris Fallin, Anjo
Vahldiek-Oberwagner, Michael LeMay, Ravi Sahita, Dean Tullsen, Deian Stefan.
IEEE Micro Top Picks, 2024

{\em Going beyond the Limits of SFI: Flexible and Secure Hardware-Assisted
In-Process Isolation with HFI.} Shravan Narayan, Tal Garfinkel, Mohammadkazem
Taram, Joey Rudek, Daniel Moghimi, Evan Johnson, Chris Fallin, Anjo Vahldiek-Oberwagner,
Michael LeMay, Ravi Sahita, Dean Tullsen, Deian Stefan.  In
Architectural Support for Programming Languages and Operating Systems (ASPLOS),
2023

{\em Segue and ColorGuard: Optimizing SFI Performance and Scalability on Modern
x86.} Shravan Narayan, Tal Garfinkel, Evan Johnson, David Thein, Joey Rudek,
Michael LeMay, Anjo Vahldiek-Oberwagner, Deian Stefan, Dean Tullsen. The 17th
Workshop on Programming Languages and Analysis for Security (PLAS), 2022

{\em The Road to Less Trusted Code: Lowering the Barrier to In-Process
Sandboxing} Tal Garfinkel, Shravan Narayan, Craig Disselkoen, Hovav Shacham,
Deian Stefan. In Usenix login; Winter 2020

{\em Retrofitting Fine Grain Isolation in the Firefox Renderer} Shravan
Narayan, Craig Disselkoen, Tal Garfinkel, Nathan Froyd, Eric Rahm, Sorin
Lerner, Hovav Shacham, Deian Stefan. In Usenix Security Symposium, 2020.

{\em Progressive Memory Safety in WebAssembly} Craig Disselkoen, John Renner,
Conrad Watt, Tal Garfinkel, Amit Levy, Deian Stefan. In Hardware and Architectural
Support for Security and Privacy, 2019.

{\em Gobi: WebAssembly as a Practical Path to Library Sandboxing} Shravan
Narayan, Tal Garfinkel, Sorin Lerner, Hovav Shacham, Deian Stefan.
arXiv:1912.02285, 2019

\pagebreak

{\em Trestle: Bridging the Performance and Safety Divide in WebAssembly}
Craig Disselkoen, Tal Garfinkel, Deian Stefan, Conrad Watt. In Principles of
Secure Compilation, 2019.


{\em Towards Practical Default-On Multi-Core Record/Replay}
Ali Mashtizadeh, Tal Garfinkel, David Terei, David Mazieres, and Mendel
Rosenblum.  In Proceedings of Symposium on Architectural Support for Programming
Languages and Operating Systems April, 2017.


{\em XvMotion: Unified Virtual Machine Migration over Long Distance} Ali
Mashtizadeh, Min Cai, Gabriel Tarasuk-Levin, Ricardo Koller, Tal Garfinkel,
Sreekanth Setty. In Usenix Annual Technical Conference, 2014.

{\em The Design and Evolution of Live Storage Migration in VMware ESX} Ali
Mashtizadeh, Emré Celebi, Tal Garfinkel, and Min Cai. In Usenix Annual Technical
Conference, 2011.

{\em Multi-stage Replay with Crosscut} Jim Chow, Dominic Lucchetti, Tal Garfinkel, Geoffrey Lefebvre, Ryan Gardner, Joshua Mason, Sam Small, and Peter M. Chen. In Conf. on Virtual Execution Environments, 2010.


{\em Virtual Machine Contracts for Datacenter and Cloud Computing Environments } Jeanna Mathews, Tal Garfinkel, Christopher Hoff, and Jeff Wheeler. In Workshop on Automated Control for Datacenters and Clouds, 2009.

{\em Towards Application Security on Untrusted Operating Systems}  Dan
Ports and Tal Garfinkel\newline
In Usenix Workshop on Hot Topics in Security, 2008.


{\em Decoupling Dynamic Program Analysis from Execution in Virtual
Environments} Jim Chow, Tal Garfinkel, and Peter M. Chen
In USENIX Annual Technical Conference, 2008.


{\em Overshadow: A Virtualization-Based Approach to Retrofitting
Protection in Commodity Operating Systems} Xiaoxin Chen, Tal
Garfinkel, E. Christopher Lewis, Pratap Subrahmanyam, Carl A.
Waldspurger, Dan Boneh, Jeffrey Dwoskin, Dan R.K. Ports.
In Architectural Support for Programming Languages and Operating
Systems, 2008.


{\em What Virtualization can do for Security} Tal Garfinkel, Andrew
Warfield ;login: The USENIX Magazine, December 2007


{\em Reducing Shoulder-surfing by Using Gaze-based Password Entry}
Manu Kumar, Tal Garfinkel, Dan Boneh, Terry Winograd. In Symposium On Usable
Privacy and Security, 2007.

{\em Compatibility is Not Transparency: VMM Detection Myths and
Realities} Tal Garfinkel, Keith Adams, Andrew Warfield, Jason
Franklin.  In Workshop on Hot Topics in Operating Systems, 2007.

{\em SANE: A Protection Architecture for Enterprise Networks} Martin
Casado,Tal Garfinkel, Aditya Akella, Michael Freedman, Dan Boneh,
Nick McKweon, Scott Shenker USENIX Security Symposium, 2006.

{\em  Virtualization Aware File Systems: Getting Beyond the Limitations
of Virtual Disks} Ben Pfaff,Tal Garfinkel, Mendel Rosenblum.  Networked Systems
Design and Implementation, 2006.

{\em Opportunistic Measurement: Extracting Insight from Spurious
Traffic} Martin Casado,Tal Garfinkel, Weidong Cu, Vern Paxson and Stefan
Savage. Workshop on Hot Topics in Networks, 2005.

{\em When Virtual is Harder than Real: Security Challenges in Virtual Machine
Based Computing Environments} Tal Garfinkel and Mendel Rosenblum 10th Workshop
on Hot Topics in Operating Systems, 2005.

{\em Shredding Your Garbage: Reducing Data Lifetime Through Secure
Deallocation} Jim Chow, Ben Pfaff, Tal Garfinkel, and Mendel Rosenblum
USENIX Security Symposium, 2005.

{\em  Virtual Machine Monitors: Current Technology and Future
Trends} Mendel Rosenblum and Tal Garfinkel. IEEE Computer, 2005

{\em Data Lifetime is a Systems Problem} Tal Garfinkel, Ben Pfaff, Jim
Chow, and Mendel Rosenblum.  SIGOPS European Workshop, 2004.

\pagebreak

{\em Understanding Data Lifetime via Whole System Simulation}
Jim Chow, Ben Pfaff, Tal Garfinkel, Kevin Christopher, and Mendel
Rosenblum. USENIX Security Symposium, 2004.


{\em Ostia: A Delegating Architecture for Secure System Call Interposition},
Tal Garfinkel, Ben Pfaff, Mendel Rosenblum.  In Proc. Network and Distributed
System Security Symposium, 2004.

{\em Traps and Pitfalls: Practical Problems in System Call Interposition
Based Security Tools}, Tal Garfinkel.  In Proc. Network and Distributed System
Security Symposium, 2003.

{\em A Virtual Machine Introspection Based Architecture for Intrusion
Detection}, Tal Garfinkel, Mendel Rosenblum.  In Proc. Network and Distributed
System Security Symposium, 2003.

{\em Flexible Operating Systems Support and Applications of Trusted Computing}, Tal Garfinkel, Mendel
Rosenblum, Dan Boneh. In Proc. Workshop on Hot Topics in Operating Systems, 2003.


{\em Terra: A Virtual Machine Based Platform for Trusted Computing}, Tal
Garfinkel, Ben Pfaff, Jim Chow, Mendel Rosenblum, Dan Boneh. In Proc. Symposium
on Operating System Principles, 2003.


\section{\sc Award Papers}

{\em Extended User Interrupts (xUI): Fast and Flexible Asynchronous
Notification without Polling.} Berk Aydogmus, Linsong Guo, Danial Zuberi, Tal
Garfinkel, Dean Tullsen, Amy Ousterhout, Kazem Taram.  In Architectural Support
for Programming Languages and Operating Systems (ASPLOS), 2025
\newline {\it Best Paper Award}

{\em Going beyond the Limits of SFI: Flexible and Secure Hardware-Assisted
In-Process Isolation with HFI.} Shravan Narayan, Tal Garfinkel, Mohammadkazem
Taram, Joey Rudek, Daniel Moghimi, Evan Johnson, Chris Fallin, Anjo Vahldiek-Oberwagner,
Michael LeMay, Ravi Sahita, Dean Tullsen, Deian Stefan.  In
Architectural Support for Programming Languages and Operating Systems (ASPLOS),
2023
\newline {\it Distinguished Paper Award}
\newline{\it Intel Hardware Security Academic Award (Honorable Mention) 2024}

{\em Retrofitting Fine Grain Isolation in the Firefox Renderer.} Shravan
Narayan, Craig Disselkoen, Tal Garfinkel, \newline Nathan Froyd, Eric Rahm, Sorin
Lerner, Hovav Shacham, Deian Stefan. In Usenix Security Symposium 2020.
\newline {\it Distinguished Paper Award}
\newline {\it CSAW Applied Research Competition (First Place) 2020}
\newline {\it NSA Cybersecurity Research Paper Competition (Honorable Mention) 2022}
\newline {\it IEEE Cybersecurity Practice Award 2022}

{\it A Virtual Machine Introspection Based Architecture for Intrusion
Detection.}\newline Tal Garfinkel and Mendel Rosenblum. Network and
Distributed System Security Symposium (NDSS), 2003.
\newline{\it Test of Time Award---NDSS 2019}
\newline{\it `` [opened] the floodgates on a tremendous amount of research
and derivative tools...one of three highest impact papers, and the most highly cited
NDSS paper, from 1995-2009.''}

{\em Decoupling Dynamic Program Analysis from Execution in Virtual \newline
Environments} Jim Chow, Tal Garfinkel, and Peter M. Chen
\newline USENIX Annual Technical Conference 2008
\newline{\it Best Paper Award}

{\em Understanding Data Lifetime via Whole System Simulation} \newline Jim
Chow, Ben Pfaff, Tal Garfinkel, Kevin Christopher, and Mendel Rosenblum
\newline{\it USENIX Security Symposium 2004, Best Paper Award}


\section{\sc Other Honors}

UC Berkeley Computer Science,  Citation for Distinction in Research, 2001 \newline
CRA Outstanding Undergraduate Award (Honorable Mention), 2001 \newline
Holmes Memorial Scholarship in Computer Science, 1997 \newline

\pagebreak 

\section{\sc Professional Activities}


{\em Program Committee}\newline
Usenix Workshop on Offensive Technologies (WOOT) (Founder and Program
Co-Chair), Usenix Security Symposium (Security), ACM Computer Security
Architecture Workshop (CSAW), Network and Distributed System Security Symposium
(NDSS), Symposiums on Usable Security and Privacy (SOUPS), ACM Workshop on
Virtual Machine Security (VMSec), ACM Cloud Computing Security Workshop (CCSW),
ACM Symposium on Cloud Computing (SOCC), Principles of Secure Compilation
(PriSC), Architecture Support for Programming Languages and Operating Systems
(ASPLOS). Recent PCs: Usenix Security 2023, ASPLOS 2023


{\em External Reviewer}\newline
IEEE Symposium on Security and Privacy, Workshop on Hot Topics in
Operating Systems (HOTOS), Usenix Annual Technical Conference,
Usenix Security Symposium, International Conference on Architectural
Support for Programming Languages and Operating Systems (ASPLOS),
Symposium on Operating Systems Design and Implementation (OSDI),
IEEE Computer, ACM Transactions on Computer Systems (TOCS),
IEEE Dependable Systems and Networks (DSN)

{\em Keynotes and Invited Talks}\newline
Qualcomm Product Security Summit 2025, 2024, Wasmcon 2024, Qcon 2024, Blackhat Briefings
2022, 2008 Conference on Detection of Intrusions, Malware, and Vulnerability
Assessment 2008. Usenix Security Symposium 2008, Annual Computer Security
Applications Conference 2007 and 2008.

{\em Selected Teaching}\newline
Instructor, ACACES Summer School 2015: {\em Virtualization and Security:
Architecture and Synthesis}\newline
Guest lecturer, Stanford: {\em Computer Security (CS155), Programming Languages
(CS 242)}


%\title{\em Technical Reviewer}
%\employer{\bf McMillan Publishing }
%\dates{Spring 99}
%\location{Berkeley,CA }
%%\begin{position}
%Pre-publication review of material on Linux and Systems administration.
%\end{position}


%\title{\em Board Member}
%\employer{\bf CALLUG }
%\dates{Fall 99 -- Present}
%\location{University of California, Berkeley}
%\begin{position}
%Leading and organizing a student group that seeks to educate others about the Linux operating system as well as offering help to those with Unix related problems.
%\end{position}

%\title{\em Teaching Assistant}
%\employer{\bf University of California, Berkeley}
%\dates{Fall 98}
%\location{University of California, Berkeley}
%\begin{position}
%Assisted Professor Michael Harrison with CS39e, a freshman seminar on
%computer security. Responsibilities included preparation of course
%materials and guest lecturing.
%\end{position}


%\title{\em Staff Writer}
%\employer{\bf Campus Point}
%\dates{Fall 97 -- Spring 98}
%\location{University of California, Santa Barbara}
%\begin{position}
%Wrote the literature column for the bi-weekly arts and entertainment paper at UCSB.
%\end{position}

\section{\sc Funding}

Cisco Gift, 2024
Improving Wasm Performance, Security, and Scalability on Modern x86 and ARM processors




\section{\sc Patents}

{\em System and Method of Manipulating Virtual Machine Recordings for High-Level Execution and Replay, 2012}

{\em Dynamic verification of validity of executable code, 2011}

{\em Method and system for recording a selected computer process for subsequent replay, 2011}

{\em Managing Latency Introduced by Virtualization, 2010}

{\em Cryptographic multi-shadowing with integrity verification, 2009}

{\em Accelerating replayed program execution to support decoupled program analysis, 2009}

{\em Decoupling dynamic program analysis from execution across heterogeneous systems, 2009}

{\em Decoupling dynamic program analysis from execution in virtual environments, 2009}

{\em Synchronous decoupled program analysis in virtual environments, 2009}

{\em Managing Latency Introduced by Virtualization, 2009}

{\em Transparent VMM-Assisted User-Mode Execution Control Transfer, 2009}

{\em Transparent Memory-Mapped Emulation of I/O Calls, 2009}




%NSF Research Internship in Science and Engineering, 1998  \newline
%Upsilon Pi Epsilon \newline

%\section{\sc Outside Interests}
%In my free time I enyjo


\end{resume}
\end{document}

